\documentclass[11pt]{article}
\usepackage[a4paper,margin=0.8in]{geometry}
\usepackage{calc}
\usepackage{array,longtable,booktabs}
\usepackage{xcolor}
\usepackage{hyperref}
\usepackage{enumitem}
\usepackage{fancyvrb}
\usepackage{fvextra}
\usepackage{framed}
\usepackage{color}
\fvset{breaklines=true, breakanywhere=true}
\DefineVerbatimEnvironment{Highlighting}{Verbatim}{breaklines,breakanywhere,commandchars=\\\{\}}
\usepackage{color}
\usepackage{fancyvrb}
\newcommand{\VerbBar}{|}
\newcommand{\VERB}{\Verb[commandchars=\\\{\}]}
\DefineVerbatimEnvironment{Highlighting}{Verbatim}{commandchars=\\\{\}}
% Add ',fontsize=\small' for more characters per line
\usepackage{framed}
\definecolor{shadecolor}{RGB}{248,248,248}
\newenvironment{Shaded}{\begin{snugshade}}{\end{snugshade}}
\newcommand{\AlertTok}[1]{\textcolor[rgb]{0.94,0.16,0.16}{#1}}
\newcommand{\AnnotationTok}[1]{\textcolor[rgb]{0.56,0.35,0.01}{\textbf{\textit{#1}}}}
\newcommand{\AttributeTok}[1]{\textcolor[rgb]{0.13,0.29,0.53}{#1}}
\newcommand{\BaseNTok}[1]{\textcolor[rgb]{0.00,0.00,0.81}{#1}}
\newcommand{\BuiltInTok}[1]{#1}
\newcommand{\CharTok}[1]{\textcolor[rgb]{0.31,0.60,0.02}{#1}}
\newcommand{\CommentTok}[1]{\textcolor[rgb]{0.56,0.35,0.01}{\textit{#1}}}
\newcommand{\CommentVarTok}[1]{\textcolor[rgb]{0.56,0.35,0.01}{\textbf{\textit{#1}}}}
\newcommand{\ConstantTok}[1]{\textcolor[rgb]{0.56,0.35,0.01}{#1}}
\newcommand{\ControlFlowTok}[1]{\textcolor[rgb]{0.13,0.29,0.53}{\textbf{#1}}}
\newcommand{\DataTypeTok}[1]{\textcolor[rgb]{0.13,0.29,0.53}{#1}}
\newcommand{\DecValTok}[1]{\textcolor[rgb]{0.00,0.00,0.81}{#1}}
\newcommand{\DocumentationTok}[1]{\textcolor[rgb]{0.56,0.35,0.01}{\textbf{\textit{#1}}}}
\newcommand{\ErrorTok}[1]{\textcolor[rgb]{0.64,0.00,0.00}{\textbf{#1}}}
\newcommand{\ExtensionTok}[1]{#1}
\newcommand{\FloatTok}[1]{\textcolor[rgb]{0.00,0.00,0.81}{#1}}
\newcommand{\FunctionTok}[1]{\textcolor[rgb]{0.13,0.29,0.53}{\textbf{#1}}}
\newcommand{\ImportTok}[1]{#1}
\newcommand{\InformationTok}[1]{\textcolor[rgb]{0.56,0.35,0.01}{\textbf{\textit{#1}}}}
\newcommand{\KeywordTok}[1]{\textcolor[rgb]{0.13,0.29,0.53}{\textbf{#1}}}
\newcommand{\NormalTok}[1]{#1}
\newcommand{\OperatorTok}[1]{\textcolor[rgb]{0.81,0.36,0.00}{\textbf{#1}}}
\newcommand{\OtherTok}[1]{\textcolor[rgb]{0.56,0.35,0.01}{#1}}
\newcommand{\PreprocessorTok}[1]{\textcolor[rgb]{0.56,0.35,0.01}{\textit{#1}}}
\newcommand{\RegionMarkerTok}[1]{#1}
\newcommand{\SpecialCharTok}[1]{\textcolor[rgb]{0.81,0.36,0.00}{\textbf{#1}}}
\newcommand{\SpecialStringTok}[1]{\textcolor[rgb]{0.31,0.60,0.02}{#1}}
\newcommand{\StringTok}[1]{\textcolor[rgb]{0.31,0.60,0.02}{#1}}
\newcommand{\VariableTok}[1]{\textcolor[rgb]{0.00,0.00,0.00}{#1}}
\newcommand{\VerbatimStringTok}[1]{\textcolor[rgb]{0.31,0.60,0.02}{#1}}
\newcommand{\WarningTok}[1]{\textcolor[rgb]{0.56,0.35,0.01}{\textbf{\textit{#1}}}}
\setlist[itemize]{leftmargin=*}
\setlist[enumerate]{leftmargin=*}
\hypersetup{colorlinks=true,linkcolor=blue,urlcolor=blue}
\providecommand{\tightlist}{%
  \setlength{\itemsep}{0pt}\setlength{\parskip}{0pt}}
\title{R Workshop -- Day 2 -- Importing, Cleaning and Preparing Data}
\author{Applied R for Statistical Teaching, Research and Reproducible Analysis}
\date{}
\begin{document}
\maketitle
\tableofcontents
\newpage
\hypertarget{r-workshop-day-2-importing-cleaning-and-preparing-data}{%
\section{R Workshop -- Day 2 -- Importing, Cleaning and Preparing
Data}\label{r-workshop-day-2-importing-cleaning-and-preparing-data}}

\textbf{Session 2 of 7 \textbar{} Duration: 3 hours} \textbf{Programme:}
Applied R for Statistical Teaching, Research and Reproducible Analysis

\hypertarget{learning-objectives}{%
\subsection{Learning objectives}\label{learning-objectives}}

By the end of this session you will be able to:

\begin{enumerate}
\def\labelenumi{\arabic{enumi}.}
\tightlist
\item
  Import a CSV file with \texttt{readr::read\_csv()} and explain why you
  keep a raw, untouched copy.
\item
  Inspect structure, names, ranges and missingness before transforming
  any data.
\item
  Write logical rules that flag invalid values without silently deleting
  rows.
\item
  Standardise inconsistent category labels.
\item
  Detect and remove duplicate records.
\item
  Produce a data-quality report, a missingness report and a short data
  dictionary.
\item
  Export a clean dataset to \texttt{Data/processed/} so later sessions
  can rely on it.
\end{enumerate}

\hypertarget{capstone-link}{%
\subsection{Capstone link}\label{capstone-link}}

Today the \textbf{anchor dataset} is introduced. From this point through
Day 7, every session works on this same dataset: a synthetic evaluation
of a national statistics-lecturer training programme. Today's task is to
choose a research question and turn the raw file into a documented,
analysis-ready dataset.

\textbf{Working research question for this workshop (you may adapt the
wording):} \emph{Does training track (Online, In-Person, Blended) and
class attendance predict improvement in lecturers' post-training
assessment scores, after accounting for their pre-training score?}

\hypertarget{segment-plan-3-hours}{%
\subsection{Segment plan (3 hours)}\label{segment-plan-3-hours}}

\begin{longtable}[]{@{}
  >{\raggedright\arraybackslash}p{(\columnwidth - 4\tabcolsep) * \real{0.3000}}
  >{\raggedleft\arraybackslash}p{(\columnwidth - 4\tabcolsep) * \real{0.4000}}
  >{\raggedright\arraybackslash}p{(\columnwidth - 4\tabcolsep) * \real{0.3000}}@{}}
\toprule\noalign{}
\begin{minipage}[b]{\linewidth}\raggedright
Segment
\end{minipage} & \begin{minipage}[b]{\linewidth}\raggedleft
Time
\end{minipage} & \begin{minipage}[b]{\linewidth}\raggedright
Purpose
\end{minipage} \\
\midrule\noalign{}
\endhead
\bottomrule\noalign{}
\endlastfoot
Opening and recap & 10 min & Revisit Day 1 baseline task, introduce the
anchor dataset \\
Concept bridge & 25 min & Why raw data stays untouched; what ``clean''
means \\
Guided coding & 60 min & Import, inspect, flag, standardise,
deduplicate \\
Participant practice & 45 min & Apply the pipeline end-to-end on the
anchor dataset \\
Interpretation and reporting & 25 min & Write the data-quality
narrative \\
Evidence log and capstone update & 15 min & Save outputs, log research
question and variables \\
\end{longtable}

\begin{center}\rule{0.5\linewidth}{0.5pt}\end{center}

\hypertarget{part-1-concept-bridge-core-workshop-activity}{%
\subsection{Part 1 -- Concept bridge (Core workshop
activity)}\label{part-1-concept-bridge-core-workshop-activity}}

\hypertarget{the-anchor-dataset}{%
\subsubsection{The anchor dataset}\label{the-anchor-dataset}}

\texttt{Data/raw/anchor\_dataset.csv} contains one row per training
participant, with deliberately realistic problems: missing values, a
handful of out-of-range values, a few duplicate records, and category
labels typed in different ways by different data-entry staff
(\texttt{"Online"}, \texttt{"online"}, \texttt{"ONLINE\ "}). This is not
a bug in the workshop materials -- it is the exercise.

\begin{longtable}[]{@{}lll@{}}
\toprule\noalign{}
Variable & Type & Description \\
\midrule\noalign{}
\endhead
\bottomrule\noalign{}
\endlastfoot
\texttt{participant\_id} & character & Unique lecturer identifier \\
\texttt{age} & numeric & Age in years \\
\texttt{gender} & character & Male / Female \\
\texttt{region} & character & Geopolitical zone \\
\texttt{institution\_type} & character & Public / Private \\
\texttt{years\_teaching} & numeric & Years of teaching experience \\
\texttt{training\_track} & character & Online / In-Person / Blended \\
\texttt{attendance} & numeric & Percentage of sessions attended
(0-100) \\
\texttt{satisfaction} & numeric & Likert scale, 1 (low) to 5 (high) \\
\texttt{pre\_score} & numeric & Pre-training assessment score (0-100) \\
\texttt{post\_score} & numeric & Post-training assessment score
(0-100) \\
\texttt{completed} & character & Yes / No -- finished the programme \\
\end{longtable}

\hypertarget{raw-versus-working-data}{%
\subsubsection{Raw versus working data}\label{raw-versus-working-data}}

\textbf{Never overwrite the raw file, and never edit it by hand.} Read
it into one object (e.g.~\texttt{raw\_lecturers}) and do all
transformation on a second object (e.g.~\texttt{clean\_lecturers}). If a
cleaning rule turns out to be wrong, you re-run the script against the
untouched raw file -- you do not try to remember what you changed.

\hypertarget{what-clean-means-here}{%
\subsubsection{What ``clean'' means here}\label{what-clean-means-here}}

Cleaning is not ``make the numbers look nicer.'' It is a documented,
reproducible set of decisions:

\begin{itemize}
\tightlist
\item
  Which values are \emph{impossible} (e.g.~attendance of 140\%) versus
  merely \emph{unusual} (e.g.~age 24, the youngest in the group)? Only
  impossible values become \texttt{NA}.
\item
  Which category labels are the \emph{same thing typed differently}
  versus genuinely different categories?
\item
  Which rows are \emph{exact duplicates} of another row, and should one
  copy be kept?
\end{itemize}

Every decision should be visible in the script, not only in your head.

\begin{center}\rule{0.5\linewidth}{0.5pt}\end{center}

\hypertarget{part-2-guided-coding-core-workshop-activity}{%
\subsection{Part 2 -- Guided coding (Core workshop
activity)}\label{part-2-guided-coding-core-workshop-activity}}

\hypertarget{import-with-readr}{%
\subsubsection{\texorpdfstring{2.1 Import with
\texttt{readr}}{2.1 Import with readr}}\label{import-with-readr}}

\begin{Shaded}
\begin{Highlighting}[]
\FunctionTok{library}\NormalTok{(tidyverse)}

\NormalTok{raw\_lecturers }\OtherTok{\textless{}{-}}\NormalTok{ readr}\SpecialCharTok{::}\FunctionTok{read\_csv}\NormalTok{(}\StringTok{"Data/raw/anchor\_dataset.csv"}\NormalTok{, }\AttributeTok{show\_col\_types =} \ConstantTok{FALSE}\NormalTok{)}
\end{Highlighting}
\end{Shaded}

\texttt{read\_csv()} (with an underscore, from \texttt{readr}) is
preferred over base \texttt{read.csv()} from today onward: it is faster
on larger files, guesses column types more transparently, and returns a
tibble that prints sensibly.

\hypertarget{inspect-before-you-touch-anything}{%
\subsubsection{2.2 Inspect before you touch
anything}\label{inspect-before-you-touch-anything}}

\begin{Shaded}
\begin{Highlighting}[]
\FunctionTok{glimpse}\NormalTok{(raw\_lecturers)}
\FunctionTok{summary}\NormalTok{(raw\_lecturers)}
\FunctionTok{nrow}\NormalTok{(raw\_lecturers)}
\FunctionTok{sum}\NormalTok{(}\FunctionTok{duplicated}\NormalTok{(raw\_lecturers))}
\FunctionTok{colSums}\NormalTok{(}\FunctionTok{is.na}\NormalTok{(raw\_lecturers))}
\end{Highlighting}
\end{Shaded}

Look specifically for: implausible minimums/maximums in
\texttt{summary()}, the duplicate count, and which columns carry the
most missingness.

\hypertarget{flag-invalid-values-do-not-delete-rows}{%
\subsubsection{2.3 Flag invalid values (do not delete
rows)}\label{flag-invalid-values-do-not-delete-rows}}

\begin{Shaded}
\begin{Highlighting}[]
\NormalTok{flagged\_lecturers }\OtherTok{\textless{}{-}}\NormalTok{ raw\_lecturers }\SpecialCharTok{\%\textgreater{}\%}
  \FunctionTok{mutate}\NormalTok{(}
    \AttributeTok{invalid\_age =} \SpecialCharTok{!}\FunctionTok{is.na}\NormalTok{(age) }\SpecialCharTok{\&}\NormalTok{ (age }\SpecialCharTok{\textless{}} \DecValTok{18} \SpecialCharTok{|}\NormalTok{ age }\SpecialCharTok{\textgreater{}} \DecValTok{100}\NormalTok{),}
    \AttributeTok{invalid\_attendance =} \SpecialCharTok{!}\FunctionTok{is.na}\NormalTok{(attendance) }\SpecialCharTok{\&}\NormalTok{ (attendance }\SpecialCharTok{\textless{}} \DecValTok{0} \SpecialCharTok{|}\NormalTok{ attendance }\SpecialCharTok{\textgreater{}} \DecValTok{100}\NormalTok{),}
    \AttributeTok{invalid\_satisfaction =} \SpecialCharTok{!}\FunctionTok{is.na}\NormalTok{(satisfaction) }\SpecialCharTok{\&} \SpecialCharTok{!}\NormalTok{satisfaction }\SpecialCharTok{\%in\%} \DecValTok{1}\SpecialCharTok{:}\DecValTok{5}\NormalTok{,}
    \AttributeTok{invalid\_pre\_score =} \SpecialCharTok{!}\FunctionTok{is.na}\NormalTok{(pre\_score) }\SpecialCharTok{\&}\NormalTok{ (pre\_score }\SpecialCharTok{\textless{}} \DecValTok{0} \SpecialCharTok{|}\NormalTok{ pre\_score }\SpecialCharTok{\textgreater{}} \DecValTok{100}\NormalTok{),}
    \AttributeTok{invalid\_post\_score =} \SpecialCharTok{!}\FunctionTok{is.na}\NormalTok{(post\_score) }\SpecialCharTok{\&}\NormalTok{ (post\_score }\SpecialCharTok{\textless{}} \DecValTok{0} \SpecialCharTok{|}\NormalTok{ post\_score }\SpecialCharTok{\textgreater{}} \DecValTok{100}\NormalTok{)}
\NormalTok{  )}

\NormalTok{flagged\_lecturers }\SpecialCharTok{\%\textgreater{}\%}
  \FunctionTok{summarise}\NormalTok{(}\FunctionTok{across}\NormalTok{(}\FunctionTok{starts\_with}\NormalTok{(}\StringTok{"invalid\_"}\NormalTok{), sum, }\AttributeTok{na.rm =} \ConstantTok{TRUE}\NormalTok{))}
\end{Highlighting}
\end{Shaded}

Flagging first means you can count and report invalid values
\emph{before} you remove them -- this count is exactly what goes into
your data-quality report.

\hypertarget{replace-flagged-values-with-na}{%
\subsubsection{\texorpdfstring{2.4 Replace flagged values with
\texttt{NA}}{2.4 Replace flagged values with NA}}\label{replace-flagged-values-with-na}}

\begin{Shaded}
\begin{Highlighting}[]
\NormalTok{clean\_lecturers }\OtherTok{\textless{}{-}}\NormalTok{ flagged\_lecturers }\SpecialCharTok{\%\textgreater{}\%}
  \FunctionTok{mutate}\NormalTok{(}
    \AttributeTok{age =} \FunctionTok{if\_else}\NormalTok{(invalid\_age, }\ConstantTok{NA\_real\_}\NormalTok{, age),}
    \AttributeTok{attendance =} \FunctionTok{if\_else}\NormalTok{(invalid\_attendance, }\ConstantTok{NA\_real\_}\NormalTok{, attendance),}
    \AttributeTok{satisfaction =} \FunctionTok{if\_else}\NormalTok{(invalid\_satisfaction, }\ConstantTok{NA\_real\_}\NormalTok{, satisfaction),}
    \AttributeTok{pre\_score =} \FunctionTok{if\_else}\NormalTok{(invalid\_pre\_score, }\ConstantTok{NA\_real\_}\NormalTok{, pre\_score),}
    \AttributeTok{post\_score =} \FunctionTok{if\_else}\NormalTok{(invalid\_post\_score, }\ConstantTok{NA\_real\_}\NormalTok{, post\_score)}
\NormalTok{  )}
\end{Highlighting}
\end{Shaded}

\hypertarget{standardise-category-labels}{%
\subsubsection{2.5 Standardise category
labels}\label{standardise-category-labels}}

\begin{Shaded}
\begin{Highlighting}[]
\NormalTok{clean\_lecturers }\OtherTok{\textless{}{-}}\NormalTok{ clean\_lecturers }\SpecialCharTok{\%\textgreater{}\%}
  \FunctionTok{mutate}\NormalTok{(}
    \AttributeTok{gender =} \FunctionTok{str\_trim}\NormalTok{(gender),}
    \AttributeTok{gender =} \FunctionTok{case\_when}\NormalTok{(}
      \FunctionTok{str\_to\_lower}\NormalTok{(gender) }\SpecialCharTok{\%in\%} \FunctionTok{c}\NormalTok{(}\StringTok{"m"}\NormalTok{, }\StringTok{"male"}\NormalTok{) }\SpecialCharTok{\textasciitilde{}} \StringTok{"Male"}\NormalTok{,}
      \FunctionTok{str\_to\_lower}\NormalTok{(gender) }\SpecialCharTok{\%in\%} \FunctionTok{c}\NormalTok{(}\StringTok{"f"}\NormalTok{, }\StringTok{"female"}\NormalTok{) }\SpecialCharTok{\textasciitilde{}} \StringTok{"Female"}\NormalTok{,}
      \ConstantTok{TRUE} \SpecialCharTok{\textasciitilde{}}\NormalTok{ gender}
\NormalTok{    ),}
    \AttributeTok{region =} \FunctionTok{str\_to\_title}\NormalTok{(}\FunctionTok{str\_trim}\NormalTok{(region)),}
    \AttributeTok{institution\_type =} \FunctionTok{str\_to\_title}\NormalTok{(}\FunctionTok{str\_trim}\NormalTok{(institution\_type)),}
    \AttributeTok{training\_track =} \FunctionTok{str\_trim}\NormalTok{(training\_track),}
    \AttributeTok{training\_track =} \FunctionTok{case\_when}\NormalTok{(}
      \FunctionTok{str\_to\_lower}\NormalTok{(training\_track) }\SpecialCharTok{==} \StringTok{"online"} \SpecialCharTok{\textasciitilde{}} \StringTok{"Online"}\NormalTok{,}
      \FunctionTok{str\_to\_lower}\NormalTok{(training\_track) }\SpecialCharTok{==} \StringTok{"blended"} \SpecialCharTok{\textasciitilde{}} \StringTok{"Blended"}\NormalTok{,}
      \FunctionTok{str\_to\_lower}\NormalTok{(training\_track) }\SpecialCharTok{\%in\%} \FunctionTok{c}\NormalTok{(}\StringTok{"in{-}person"}\NormalTok{, }\StringTok{"in person"}\NormalTok{, }\StringTok{"inperson"}\NormalTok{) }\SpecialCharTok{\textasciitilde{}} \StringTok{"In{-}Person"}\NormalTok{,}
      \ConstantTok{TRUE} \SpecialCharTok{\textasciitilde{}}\NormalTok{ training\_track}
\NormalTok{    )}
\NormalTok{  )}

\FunctionTok{count}\NormalTok{(clean\_lecturers, gender)}
\FunctionTok{count}\NormalTok{(clean\_lecturers, training\_track)}
\end{Highlighting}
\end{Shaded}

\texttt{str\_trim()} removes leading/trailing whitespace;
\texttt{str\_to\_lower()}/\texttt{str\_to\_title()} normalise case
before you compare or display values. \texttt{case\_when()} reads top to
bottom and stops at the first match -- always finish with a fallback
(\texttt{TRUE\ \textasciitilde{}\ ...}) so unexpected values are not
silently dropped.

\hypertarget{duplicates}{%
\subsubsection{2.6 Duplicates}\label{duplicates}}

\begin{Shaded}
\begin{Highlighting}[]
\NormalTok{n\_before }\OtherTok{\textless{}{-}} \FunctionTok{nrow}\NormalTok{(clean\_lecturers)}
\NormalTok{clean\_lecturers }\OtherTok{\textless{}{-}}\NormalTok{ clean\_lecturers }\SpecialCharTok{\%\textgreater{}\%} \FunctionTok{distinct}\NormalTok{(participant\_id, }\AttributeTok{.keep\_all =} \ConstantTok{TRUE}\NormalTok{)}
\NormalTok{n\_duplicates\_removed }\OtherTok{\textless{}{-}}\NormalTok{ n\_before }\SpecialCharTok{{-}} \FunctionTok{nrow}\NormalTok{(clean\_lecturers)}
\NormalTok{n\_duplicates\_removed}
\end{Highlighting}
\end{Shaded}

Deduplicate on the identifier \emph{after} standardising labels --
otherwise \texttt{"Online"} and \texttt{"online"} rows for the same
participant might look like different records.

\hypertarget{data-quality-and-missingness-reports}{%
\subsubsection{2.7 Data-quality and missingness
reports}\label{data-quality-and-missingness-reports}}

\begin{Shaded}
\begin{Highlighting}[]
\NormalTok{quality\_report }\OtherTok{\textless{}{-}}\NormalTok{ flagged\_lecturers }\SpecialCharTok{\%\textgreater{}\%}
  \FunctionTok{summarise}\NormalTok{(}
    \AttributeTok{rows =} \FunctionTok{n}\NormalTok{(),}
    \AttributeTok{duplicates\_removed =}\NormalTok{ n\_duplicates\_removed,}
    \AttributeTok{invalid\_age =} \FunctionTok{sum}\NormalTok{(invalid\_age, }\AttributeTok{na.rm =} \ConstantTok{TRUE}\NormalTok{),}
    \AttributeTok{invalid\_attendance =} \FunctionTok{sum}\NormalTok{(invalid\_attendance, }\AttributeTok{na.rm =} \ConstantTok{TRUE}\NormalTok{),}
    \AttributeTok{invalid\_satisfaction =} \FunctionTok{sum}\NormalTok{(invalid\_satisfaction, }\AttributeTok{na.rm =} \ConstantTok{TRUE}\NormalTok{),}
    \AttributeTok{invalid\_pre\_score =} \FunctionTok{sum}\NormalTok{(invalid\_pre\_score, }\AttributeTok{na.rm =} \ConstantTok{TRUE}\NormalTok{),}
    \AttributeTok{invalid\_post\_score =} \FunctionTok{sum}\NormalTok{(invalid\_post\_score, }\AttributeTok{na.rm =} \ConstantTok{TRUE}\NormalTok{)}
\NormalTok{  )}

\NormalTok{missingness\_report }\OtherTok{\textless{}{-}}\NormalTok{ clean\_lecturers }\SpecialCharTok{\%\textgreater{}\%}
  \FunctionTok{select}\NormalTok{(}\SpecialCharTok{{-}}\FunctionTok{starts\_with}\NormalTok{(}\StringTok{"invalid\_"}\NormalTok{)) }\SpecialCharTok{\%\textgreater{}\%}
  \FunctionTok{summarise}\NormalTok{(}\FunctionTok{across}\NormalTok{(}\FunctionTok{everything}\NormalTok{(), }\SpecialCharTok{\textasciitilde{}} \FunctionTok{sum}\NormalTok{(}\FunctionTok{is.na}\NormalTok{(.x)))) }\SpecialCharTok{\%\textgreater{}\%}
  \FunctionTok{pivot\_longer}\NormalTok{(}\FunctionTok{everything}\NormalTok{(), }\AttributeTok{names\_to =} \StringTok{"variable"}\NormalTok{, }\AttributeTok{values\_to =} \StringTok{"missing\_count"}\NormalTok{)}
\end{Highlighting}
\end{Shaded}

\hypertarget{export}{%
\subsubsection{2.8 Export}\label{export}}

\begin{Shaded}
\begin{Highlighting}[]
\NormalTok{clean\_lecturers }\OtherTok{\textless{}{-}}\NormalTok{ clean\_lecturers }\SpecialCharTok{\%\textgreater{}\%} \FunctionTok{select}\NormalTok{(}\SpecialCharTok{{-}}\FunctionTok{starts\_with}\NormalTok{(}\StringTok{"invalid\_"}\NormalTok{))}

\FunctionTok{dir.create}\NormalTok{(}\StringTok{"Data/processed"}\NormalTok{, }\AttributeTok{showWarnings =} \ConstantTok{FALSE}\NormalTok{, }\AttributeTok{recursive =} \ConstantTok{TRUE}\NormalTok{)}
\FunctionTok{write\_csv}\NormalTok{(clean\_lecturers, }\StringTok{"Data/processed/anchor\_dataset\_clean.csv"}\NormalTok{)}
\FunctionTok{write\_csv}\NormalTok{(quality\_report, }\StringTok{"Data/processed/data\_quality\_report.csv"}\NormalTok{)}
\FunctionTok{write\_csv}\NormalTok{(missingness\_report, }\StringTok{"Data/processed/missingness\_report.csv"}\NormalTok{)}
\end{Highlighting}
\end{Shaded}

\begin{center}\rule{0.5\linewidth}{0.5pt}\end{center}

\hypertarget{part-3-participant-practice-core-workshop-activity}{%
\subsection{Part 3 -- Participant practice (Core workshop
activity)}\label{part-3-participant-practice-core-workshop-activity}}

Open \texttt{Starter-Scripts/day2\_anchor\_cleaning.R}. It already
contains the skeleton above with \texttt{TODO} markers. Work through
each \texttt{TODO} yourself rather than copying Part 2 verbatim -- in
particular:

\begin{enumerate}
\def\labelenumi{\arabic{enumi}.}
\tightlist
\item
  Run the import and inspection steps and write down, in a comment,
  which columns concerned you most before cleaning.
\item
  Complete the invalid-value replacement step.
\item
  Complete the category-standardisation step, then check
  \texttt{count()} on each categorical column until every label set
  looks correct.
\item
  Remove duplicates and confirm the row count drop matches what you
  expected from \texttt{sum(duplicated(...))}.
\item
  Build the two reports and the export.
\item
  Create \texttt{Data/processed/data\_dictionary.csv} with one row per
  variable: \texttt{variable}, \texttt{type}, \texttt{description},
  \texttt{valid\_range\_or\_levels}.
\end{enumerate}

A facilitator-reviewed solution is in
\texttt{Solution-Scripts/day2\_anchor\_cleaning\_solution.R}.

\begin{center}\rule{0.5\linewidth}{0.5pt}\end{center}

\hypertarget{optional-demonstration}{%
\subsection{Optional demonstration}\label{optional-demonstration}}

\begin{itemize}
\tightlist
\item
  \textbf{Reshaping} with
  \texttt{pivot\_longer()}/\texttt{pivot\_wider()} -- useful once you
  start comparing pre- and post-scores side by side as one ``time''
  column rather than two columns.
\item
  \textbf{Importing other formats:} \texttt{readxl::read\_excel()},
  \texttt{haven::read\_sav()} (SPSS), \texttt{haven::read\_dta()}
  (Stata), \texttt{haven::read\_sas()} (SAS). Demonstrate one if
  participants' own data comes from these formats.
\end{itemize}

\begin{Shaded}
\begin{Highlighting}[]
\CommentTok{\# Optional: reshape pre/post scores into long format for later plotting}
\NormalTok{clean\_lecturers }\SpecialCharTok{\%\textgreater{}\%}
  \FunctionTok{select}\NormalTok{(participant\_id, pre\_score, post\_score) }\SpecialCharTok{\%\textgreater{}\%}
  \FunctionTok{pivot\_longer}\NormalTok{(}\AttributeTok{cols =} \FunctionTok{c}\NormalTok{(pre\_score, post\_score),}
               \AttributeTok{names\_to =} \StringTok{"timing"}\NormalTok{, }\AttributeTok{values\_to =} \StringTok{"score"}\NormalTok{)}
\end{Highlighting}
\end{Shaded}

\begin{center}\rule{0.5\linewidth}{0.5pt}\end{center}

\hypertarget{reference-or-take-home-material}{%
\subsection{Reference or take-home
material}\label{reference-or-take-home-material}}

\begin{itemize}
\tightlist
\item
  Keeping an audit trail: instead of overwriting values, some teams add
  a \texttt{*\_original} column before correcting a value, so the
  original is always recoverable.
\item
  More validation-rule examples (date ranges, cross-field checks such as
  \texttt{post\_score} requiring a non-missing \texttt{pre\_score}).
\end{itemize}

\begin{center}\rule{0.5\linewidth}{0.5pt}\end{center}

\hypertarget{daily-deliverable}{%
\subsection{Daily deliverable}\label{daily-deliverable}}

A clean anchor dataset, a data-quality report, a missingness report, and
a short data dictionary, all saved in \texttt{Data/processed/}, plus a
one-paragraph statement of your chosen research question.

\hypertarget{evidence-log-entry}{%
\subsubsection{Evidence log entry}\label{evidence-log-entry}}

Complete the Day 2 row in \texttt{Workshop-Evidence-Log.md}: note your
research question, confirm the clean dataset and dictionary are saved,
and record one cleaning decision you are unsure about so you can revisit
it later.

\begin{center}\rule{0.5\linewidth}{0.5pt}\end{center}

\hypertarget{facilitator-notes}{%
\subsection{Facilitator notes}\label{facilitator-notes}}

\begin{itemize}
\tightlist
\item
  The dataset is generated with a fixed random seed, so every
  participant's raw file is identical -- this lets you state exact
  expected counts (240 unique participants, 6 duplicate rows, around 5\%
  missingness per column) if you want to use them as a checkpoint.
\item
  Watch for participants who delete invalid rows outright instead of
  flagging then setting to \texttt{NA}. Deleting rows silently changes
  the denominator for every later analysis; this is the single most
  common and most consequential mistake on Day 2.
\item
  If a participant's \texttt{training\_track} recode misses a label,
  \texttt{count(clean\_lecturers,\ training\_track)} will show it
  immediately -- use this as a live debugging demonstration.
\end{itemize}
\end{document}
